\documentclass[11pt]{article}
\usepackage[margin=1in]{geometry}
\usepackage{amsmath,amssymb}
\usepackage{booktabs}
\usepackage{enumitem}
\usepackage{microtype}
\usepackage[hidelinks]{hyperref}
\usepackage{titlesec}
\titleformat{\section}{\normalfont\large\bfseries}{\thesection.}{0.6em}{}
\titlespacing*{\section}{0pt}{1.4em}{0.5em}
\newcommand{\hdr}[1]{\medskip\noindent\textbf{#1}\par\nopagebreak}
\setlist{nosep,leftmargin=1.5em,topsep=2pt}
\frenchspacing

\title{\vspace{-1.2cm}\bfseries Terrain-Aware Semantic Segmentation for Autonomous Mars
Rover Driving: CNNs vs.\ Transformers vs.\ Foundation Models}
\author{}
\date{}

\begin{document}
\maketitle
\vspace{-1.3cm}
\noindent\textbf{Course:} EN.705.742 --- Advanced Applied Machine Learning, Summer 2026
\quad\textbf{Author:} John Roth \quad\textbf{Submission:} Module 2 Research Topic
\par\smallskip\noindent\rule{0.45\textwidth}{0.4pt}

\section{Introduction and Goals}

\hdr{Research Topic}
This project investigates which family of modern semantic-segmentation
models---convolutional, transformer, or foundation---best classifies Martian terrain from rover
camera imagery, the core perception primitive for \emph{autonomous drivability} assessment. A
planetary rover that can label each pixel of the scene ahead as \emph{soil}, \emph{bedrock},
\emph{sand}, or \emph{big rock} can plan longer, safer autonomous drives without waiting for a
round-trip command from Earth. The stakes are concrete: Curiosity sustained significant wheel damage
driving over sharp embedded bedrock, and onboard terrain classification (NASA JPL's SPOC system
[14]) was developed precisely to mitigate such hazards. The NASA \textbf{AI4Mars} dataset [1,18]
---\,$\sim$326K crowdsourced segmentation labels over $\sim$35K images from the Curiosity,
Perseverance, Opportunity, and Spirit rovers\,---\,makes a rigorous, data-driven study of this
problem possible for the first time.

\hdr{A central architectural question motivates the comparison}
A Martian scene has structure at two scales: fine local texture and boundaries (the silhouette of a
small rock, the lip of an outcrop) and large homogeneous regions (broad soil flats, dune fields).
Convolutional encoder--decoders (U-Net [2], DeepLabV3+ [3]) bake in strong local inductive bias and
multi-scale context through pooling and atrous convolution; the SegFormer transformer [4] models
\emph{global} context directly through self-attention over the whole image; and vision foundation
models (SAM [5], DINOv3 [6]) bring web- and \emph{satellite}-scale pretrained priors. These three
inductive biases over the \emph{same} pixel-labeling task make AI4Mars a well-motivated testbed for
the question of which prior best captures the structure of planetary terrain.

\hdr{Design for a robust, fair comparison}
Terrain class is \emph{directly observable in the pixels}, so a meaningful positive result is
expected; the scientific value lies in \emph{which} model wins and \emph{why}. Five deliberate design
choices make the comparison fair and the conclusions defensible:
\begin{enumerate}
  \item \textbf{Controlled architecture comparison.} Models share one training recipe (optimizer,
  schedule, augmentation, crop size, loss), so the architecture is the only variable.
  \item \textbf{Transfer as a controlled lever.} Each CNN is run with an ImageNet-pretrained encoder
  \emph{and} from an identical random initialization, isolating the value of transfer (H2).
  \item \textbf{Honest evaluation.} Models are selected on a validation split and reported on the
  official \emph{expert-labeled} AI4Mars test set; a separate \emph{cross-rover} test (train on
  Curiosity, test on Opportunity/Spirit) measures real generalization (H4).
  \item \textbf{Imbalance-aware metrics.} Soil dominates the pixel distribution, so a class-weighted
  loss is used and \emph{per-class} IoU is reported, not just the mean.
  \item \textbf{Pre-registered significance.} Model differences are claimed only when significant
  under a paired bootstrap (and McNemar's test) with Holm correction.
\end{enumerate}

\hdr{Pipeline Overview}
\begin{verbatim}
NASA AI4Mars (Curiosity/Perseverance/Opportunity/Spirit)
    grayscale Navcam images  +  crowdsourced labels (0-3; 255 = ignore)
        |
SegDataset: ignore-mask handling, by-image splits, class-weighted loss
        |
   +----+--------------------+--------------------+-----------------+
   v                         v                    v                 v
Baseline U-Net      U-Net / DeepLabV3+        SegFormer        Foundation
(from scratch)      (ImageNet vs scratch)    (MiT encoder)    (SAM / DINOv2)
   |                         |                    |                 |
   +----------> CE + Dice loss (ignore_index), AdamW, AMP on GPU <--+
        |
Evaluate: mIoU / per-class IoU / pixel-acc / boundary-F1
        |
Paired bootstrap + McNemar (Holm)  ->  per-class & cross-rover stratified analysis
\end{verbatim}

\hdr{Project Goals}
\begin{itemize}
  \item Build a reproducible, leakage-safe AI4Mars segmentation pipeline (by-image splits;
  ignore-mask exclusion from loss and metrics).
  \item Compare CNN (U-Net, DeepLabV3+) vs.\ transformer (SegFormer) vs.\ foundation (SAM/DINOv2)
  architectures under one training recipe.
  \item Quantify the value of ImageNet transfer learning vs.\ from-scratch training.
  \item Test cross-rover generalization (Curiosity $\rightarrow$ Opportunity/Spirit) and whether
  transfer/augmentation narrows the domain gap.
  \item Determine which inductive bias---local convolution, global attention, or foundation
  prior---best captures Martian terrain, and report drivability-relevant per-class accuracy.
\end{itemize}

\hdr{ML Problem Statement}
This is a \textbf{dense multi-class semantic segmentation} problem. Given a grayscale rover image
$\mathbf{x}\in\mathbb{R}^{H\times W}$ (replicated to three channels for pretrained encoders), a model
$f_\theta:\mathbb{R}^{H\times W\times 3}\!\to\!\mathbb{R}^{H\times W\times C}$ produces per-pixel
class logits over $C=4$ terrain classes, yielding a label map
$\hat{\mathbf{y}}\in\{0,\dots,3\}^{H\times W}$. Training minimizes a class-weighted cross-entropy plus
Dice loss computed \emph{only over valid pixels} (the ignore label $255$---unlabeled, rover-self, and
$>$30\,m range---is excluded). The headline metric is mean Intersection-over-Union (mIoU). The
applied ML methods are U-Net, DeepLabV3+, and SegFormer, with ImageNet/foundation transfer learning,
compared against a from-scratch baseline.

\section{Problem, Candidate Solution, and Research Hypotheses}

\hdr{Problem Statement}
Autonomous driving is the bottleneck on planetary-rover science return: every meter driven under
ground-in-the-loop control costs hours of communication latency. Reliable onboard terrain
segmentation enables longer autonomous traverses and hazard avoidance. AI4Mars [1] supplies the
labeled imagery to learn this, and prior work has reported strong single-model results
(SPOC [14]; MarsSeg [13]). What is \emph{not} established is a controlled, significance-tested
comparison across the modern architecture families, paired with an honest measurement of how a model
trained on one rover generalizes to another camera and terrain distribution.

\smallskip\noindent The gap this project fills is: \textbf{which inductive bias---local convolution,
global self-attention, or a pretrained foundation prior---best segments Martian terrain; how much
does transfer learning help; and does a model trained on Curiosity generalize to the Opportunity and
Spirit rovers?}

\hdr{Null Hypothesis (H0)}
No deep segmentation model significantly outperforms a simple from-scratch baseline on mean IoU,
implying that architecture and transfer choices do not matter for this task.

\hdr{Candidate Solution}
\begin{enumerate}
  \item Download and verify the open AI4Mars merged release; build a \texttt{SegDataset} that pairs
  each image with its label map, excludes the ignore region, and splits \emph{by image}.
  \item Train a from-scratch baseline U-Net to establish the H0/H1 yardstick.
  \item Train U-Net and DeepLabV3+ with ResNet/EfficientNet encoders, each \emph{with} ImageNet
  pretraining and \emph{without} (from scratch), under one recipe (class-weighted CE\,+\,Dice, AdamW,
  cosine schedule, early stopping on val mIoU).
  \item Train SegFormer (hierarchical MiT encoder) as the transformer arm.
  \item Evaluate a foundation reference: SAM zero-shot mask proposals mapped to terrain classes,
  and/or a DINOv2 backbone with a trained segmentation head.
  \item Evaluate on the expert test set and the cross-rover (MER) set; compare all models with a
  pre-registered paired bootstrap / McNemar test under Holm correction.
\end{enumerate}

\hdr{Research Hypotheses}
\begin{itemize}
  \item \textbf{H1 (Signal Existence):} At least one deep model (U-Net, DeepLabV3+, or SegFormer)
  significantly outperforms the from-scratch baseline on mIoU (paired-bootstrap significant),
  confirming that learned multi-scale representations matter for Mars terrain.
  \item \textbf{H2 (Transfer Learning):} A model with an ImageNet-pretrained encoder significantly
  outperforms its \emph{identically configured} from-scratch counterpart, showing that natural-image
  priors transfer to grayscale planetary imagery despite the large domain gap.
  \item \textbf{H3 (Architecture / Inductive Bias):} SegFormer's global attention yields higher IoU
  on large homogeneous classes (soil, sand), while CNN decoders yield higher IoU on the small,
  boundary-dominated \emph{big rock} class---i.e.\ the winner depends on the class, tested via
  per-class IoU strata.
  \item \textbf{H4 (Cross-Rover Robustness):} A model trained on Curiosity generalizes to the
  Opportunity/Spirit (MER) imagery with a bounded mIoU drop, and ImageNet pretraining plus
  augmentation measurably narrow the cross-rover gap relative to from-scratch training.
  \item \textbf{H5 (Foundation-Model Reference):} A vision foundation model---SAM used zero-shot, or a
  \textbf{DINOv3 backbone pretrained on Earth \emph{satellite} imagery} (SAT-493M) with a light
  trained head---matches or exceeds the from-scratch baseline. This tests a striking cross-planet,
  cross-viewpoint transfer question: do self-supervised features learned from Earth orbital RGB
  generalize to Mars ground-level grayscale terrain?
\end{itemize}

\section{Applied Machine Learning Method and Dataset}

\hdr{Dataset}
\begin{center}
\begin{tabular}{@{}lll@{}}
\toprule
Source & Description & Access \\
\midrule
NASA AI4Mars [1,18] & $\sim$326K crowdsourced segmentation labels over & Free, open \\
(merged-0.6) & $\sim$35K rover images; 4 terrain classes (0--3); & (Zenodo) \\
 & 255 = NULL/ignore (unlabeled, rover, $>$30\,m) & \\
AI4Mars MSL & Curiosity Navcam: $\sim$16{,}064 train + 322 expert & Free, open \\
(Curiosity) & test images (\texttt{masked-gold-min$k$-100agree}) & \\
AI4Mars MER & Opportunity / Spirit imagery for the & Free, open \\
(cross-rover) & cross-rover generalization test (H4) & \\
ImageNet-1k [11] & 1.28M natural images; CNN encoder pretraining (H2) & Public \\
SAT-493M [6] & 493M Maxar RGB orthoimages (0.6\,m); DINOv3 & Public \\
 & satellite backbone pretraining (H5) & \\
\bottomrule
\end{tabular}
\end{center}

\hdr{Protocol and leakage controls}
Splits are made \emph{by image} so no frame appears in two splits; the official expert test set is
held out for final reporting. The ignore label (255) is excluded from both the loss and every metric.
Because soil dominates the pixel distribution, the loss is inverse-frequency class-weighted and
per-class IoU is always reported. Augmentation preserves the scene geometry: horizontal flip and mild
photometric jitter only (no vertical flip, which would invert the horizon). Grayscale inputs are
replicated to three channels and ImageNet-normalized so pretrained encoders see a compatible
distribution. Every run records its seed, git SHA, configuration hash, and resolved package versions
in a manifest for reproducibility.

\hdr{ML Methods}

\noindent\textbf{Baseline: from-scratch U-Net.} A small U-Net trained from random initialization on
AI4Mars only---the H0/H1 yardstick that isolates the value of capacity and transfer.

\smallskip\noindent\textbf{Model 1: U-Net (CNN encoder--decoder) [2].} A contracting encoder and
expanding decoder joined by \emph{skip connections} that fuse high-resolution local detail with
deep semantic context. With encoder feature map $\mathbf{e}_\ell$ and decoder map
$\mathbf{d}_{\ell+1}$, each decoder stage computes
\[
\mathbf{d}_\ell = \mathrm{Conv}\big(\,[\,\mathrm{Up}(\mathbf{d}_{\ell+1})\,;\,\mathbf{e}_\ell\,]\,\big),
\]
the concatenation $[\cdot\,;\cdot]$ being the mechanism that recovers sharp boundaries. Run with
\textbf{ResNet-34} and \textbf{EfficientNet-B0} encoders (\texttt{segmentation-models-pytorch}), each
ImageNet-1k-pretrained vs.\ from scratch.

\smallskip\noindent\textbf{Model 2: DeepLabV3+ [3].} Adds an Atrous Spatial Pyramid Pooling module
that probes context at multiple dilation rates $r$ without losing resolution. For input $x$ and
filter $w$, atrous convolution is
\[
y[i] = \sum_{k} x[i + r\,k]\,w[k],
\]
so larger $r$ enlarges the receptive field at constant cost---well suited to broad homogeneous Martian
terrain. Run with a \textbf{ResNet-50} encoder, ImageNet-1k-pretrained vs.\ from scratch.

\smallskip\noindent\textbf{Model 3: SegFormer (transformer) [4].} A hierarchical Mix-Transformer
encoder (\textbf{MiT-B0} and \textbf{MiT-B2}, via \texttt{transformers}) with efficient self-attention
and a lightweight all-MLP decoder. Self-attention models \emph{global} pixel context:
\[
\mathrm{Attention}(Q,K,V) = \mathrm{softmax}\!\Big(\tfrac{Q K^{\top}}{\sqrt{d}}\Big)V,
\]
with a sequence-reduction trick that makes attention tractable at image resolution. This is the
primary test of whether global attention beats convolutional locality on planetary terrain.

\smallskip\noindent\textbf{Foundation reference [5,6].} SAM [5] (\texttt{sam\_vit\_b}) is a promptable
segmentation model pretrained on $>$1B masks; we evaluate its zero-shot mask proposals mapped to
terrain classes. \textbf{DINOv3} [6] provides self-supervised ViT features [8]; we use the
\textbf{ViT-L/16} variant pretrained on \textbf{SAT-493M}
(\texttt{facebook/dinov3-vitl16-pretrain-sat493m}; 493M Maxar RGB orthoimages at 0.6\,m) as a
\emph{frozen} backbone with a light trained decoder head---a domain-relevant prior (terrain and
geologic texture seen from orbit) far closer to Mars than ImageNet natural images. Both probe how far
large-scale visual priors transfer to Martian terrain (H5).

\smallskip\noindent\textbf{Loss.} All trained models minimize a combined objective on valid pixels,
\[
\mathcal{L} = \mathcal{L}_{\mathrm{CE}}^{\,w}(\hat{y}, y) + \lambda\,\big(1 - \mathrm{Dice}(\hat{y}, y)\big),
\qquad
\mathrm{Dice} = \frac{2\sum p\,g}{\sum p + \sum g},
\]
where $\mathcal{L}_{\mathrm{CE}}^{\,w}$ is inverse-frequency-weighted cross-entropy (ignore index
255), $p$ are predicted soft probabilities, $g$ the one-hot targets, and $\lambda$ balances the
terms. Dice [9] directly targets the IoU-like overlap and counteracts class imbalance.

\hdr{Architecture Comparison Summary}
\begin{center}
\begin{tabular}{@{}llll@{}}
\toprule
Model & Inductive bias & Context mechanism & Key hyperparameters \\
\midrule
U-Net [2] & Local convolution & Skip connections & ResNet-34 / EfficientNet-B0 \\
DeepLabV3+ [3] & Local + multi-scale & Atrous pyramid (ASPP) & ResNet-50, dilation rates \\
SegFormer [4] & Global attention & Efficient self-attention & MiT-B0 / MiT-B2, decoder dim \\
SAM / DINOv3-SAT [5,6] & Foundation prior & Pretrained ViT [8] & ViT-L/16-sat493m; SAM ViT-B \\
\bottomrule
\end{tabular}
\end{center}

\hdr{Hyperparameter Search Plan}
\begin{center}
\begin{tabular}{@{}lll@{}}
\toprule
Hyperparameter & Search range & Method \\
\midrule
Encoder & ResNet-34, EfficientNet-B0, MiT-B0/B2 & Grid \\
Pretraining & ImageNet, none (scratch) & Grid (H2) \\
Learning rate & 1e-4, 3e-4, 6e-4 & Grid \\
Crop / input size & 384, 512 & Grid \\
Batch size & 8, 16 & Grid \\
Dice weight $\lambda$ & 0.5, 1.0 & Grid \\
\bottomrule
\end{tabular}
\end{center}
\noindent All selection is performed on a held-out validation split; the expert test set is scored
only for the best validated configuration per model.

\hdr{Performance Metrics}
\begin{itemize}
  \item \textbf{mean IoU} (primary) and \textbf{per-class IoU} for soil, bedrock, sand, big rock
  (per-class drives H3 and is the drivability-relevant view).
  \item \textbf{Pixel accuracy} and \textbf{boundary-F1} (boundary quality matters for hazard edges).
  \item \textbf{Statistical significance:} a \emph{paired bootstrap} over test images and
  \emph{McNemar's test} [15] on per-pixel correctness for pairwise model comparison, Holm-corrected
  within each hypothesis family, so differences are reported as wins only when significant.
  \item \textbf{Stratified analysis:} per-class IoU (tests H3); in-rover vs.\ cross-rover (tests H4);
  ImageNet-pretrained vs.\ from-scratch (tests H2).
\end{itemize}

\hdr{Implementation Stack}
\begin{itemize}
  \item \texttt{PyTorch} / \texttt{torchvision} --- model implementation and training.
  \item \texttt{segmentation-models-pytorch} --- U-Net / DeepLabV3+ with ImageNet-pretrained encoders.
  \item \texttt{transformers} --- SegFormer (MiT) segmentation; \texttt{segment-anything}, \texttt{timm}
  (DINOv2) --- foundation references.
  \item \texttt{albumentations}, \texttt{opencv}, \texttt{scikit-image} --- image IO, augmentation,
  boundary metrics.
  \item \texttt{scikit-learn} / \texttt{scipy} --- bootstrap CIs, McNemar test, multiple-comparison
  correction; \texttt{matplotlib} --- segmentation-overlay figures.
\end{itemize}

\section{References}
\begin{enumerate}[leftmargin=2.2em,label={[\arabic*]}]
  \item Swan, R.~M., Atha, D., Leopold, H.~A., et al. (2021). \emph{AI4MARS: A Dataset for
  Terrain-Aware Autonomous Driving on Mars}. CVPR Workshops 2021.
  \item Ronneberger, O., Fischer, P., \& Brox, T. (2015). \emph{U-Net: Convolutional Networks for
  Biomedical Image Segmentation}. MICCAI 2015. arXiv:1505.04597.
  \item Chen, L.-C., Zhu, Y., Papandreou, G., Schroff, F., \& Adam, H. (2018). \emph{Encoder-Decoder
  with Atrous Separable Convolution for Semantic Image Segmentation (DeepLabV3+)}. ECCV 2018.
  arXiv:1802.02611.
  \item Xie, E., Wang, W., Yu, Z., Anandkumar, A., Alvarez, J.~M., \& Luo, P. (2021). \emph{SegFormer:
  Simple and Efficient Design for Semantic Segmentation with Transformers}. NeurIPS 2021.
  arXiv:2105.15203.
  \item Kirillov, A., Mintun, E., Ravi, N., et al. (2023). \emph{Segment Anything}. ICCV 2023.
  arXiv:2304.02643.
  \item \emph{DINOv3: Self-supervised Learning for Vision at Scale} (2025). Meta AI Research.
  arXiv:2508.10104. Satellite variant (ViT pretrained on SAT-493M Maxar imagery):
  \url{https://huggingface.co/facebook/dinov3-vit7b16-pretrain-sat493m}; code:
  \url{https://github.com/facebookresearch/dinov3}.
  \item He, K., Zhang, X., Ren, S., \& Sun, J. (2016). \emph{Deep Residual Learning for Image
  Recognition (ResNet)}. CVPR 2016. arXiv:1512.03385.
  \item Dosovitskiy, A., et al. (2021). \emph{An Image is Worth 16$\times$16 Words: Transformers for
  Image Recognition at Scale (ViT)}. ICLR 2021. arXiv:2010.11929.
  \item Milletari, F., Navab, N., \& Ahmadi, S.-A. (2016). \emph{V-Net: Fully Convolutional Neural
  Networks for Volumetric Medical Image Segmentation (Dice loss)}. 3DV 2016. arXiv:1606.04797.
  \item Long, J., Shelhamer, E., \& Darrell, T. (2015). \emph{Fully Convolutional Networks for
  Semantic Segmentation}. CVPR 2015. arXiv:1411.4038.
  \item Deng, J., Dong, W., Socher, R., Li, L.-J., Li, K., \& Fei-Fei, L. (2009). \emph{ImageNet: A
  Large-Scale Hierarchical Image Database}. CVPR 2009.
  \item Loshchilov, I., \& Hutter, F. (2019). \emph{Decoupled Weight Decay Regularization (AdamW)}.
  ICLR 2019. arXiv:1711.05101.
  \item \emph{MarsSeg: Mars Surface Semantic Segmentation with Multi-level Extractor and Connector}
  (2024). arXiv:2404.04155.
  \item Rothrock, B., Kennedy, R., Cunningham, C., Papon, J., Heverly, M., \& Ono, M. (2016).
  \emph{SPOC: Deep Learning-based Terrain Classification for Mars Rover Missions}. AIAA SPACE 2016.
  \item Dietterich, T.~G. (1998). \emph{Approximate Statistical Tests for Comparing Supervised
  Classification Learning Algorithms (McNemar)}. Neural Computation, 10(7), 1895--1923.
  \item Efron, B., \& Tibshirani, R.~J. (1993). \emph{An Introduction to the Bootstrap}. Chapman \&
  Hall.
  \item Everingham, M., Van Gool, L., Williams, C.~K.~I., Winn, J., \& Zisserman, A. (2010).
  \emph{The PASCAL Visual Object Classes (VOC) Challenge}. IJCV, 88(2), 303--338.
  \item NASA. \emph{AI4MARS: A Dataset for Terrain-Aware Autonomous Driving on Mars} (merged-0.6).
  Zenodo. \url{https://zenodo.org/records/15995036}.
\end{enumerate}

\end{document}
